\reviewsection{Summary of the 10-Minute Lesson Taught}

The lesson was titled \textit{Many Ways to Be a Good Classmate}. The student-friendly goal was: \textit{I can name two ways to include and support a classmate who communicates, learns, moves, or experiences the room differently.}

I began with a visible first/next/last sequence. First, students learned one respectful idea about autism. Next, they practiced with a familiar classroom situation. Last, they chose one supportive action they could use. This predictable structure reflected my placement observations that young learners were more successful when directions were short, sequenced, visible, and connected to an immediate task.

\begin{figure}[htbp]
\centering
\resizebox{0.96\textwidth}{!}{\begin{tikzpicture}[step/.style={rounded corners=3mm,draw=DeepBlue,very thick,minimum width=2.35in,minimum height=1in,align=center},arrow/.style={-{Latex[length=3mm]},very thick,draw=PrimaryRed}]
\node[step,fill=SoftBlue] (a) {\textbf{0--2 minutes}\\Preview first/next/last\\Introduce autism respectfully};
\node[step,fill=SoftGreen,right=7mm of a] (b) {\textbf{2--4 minutes}\\Normalize communication\\and sensory differences};
\node[step,fill=SoftGold,right=7mm of b] (c) {\textbf{4--8 minutes}\\Teach five peer actions\\Practice a coding scenario};
\node[step,fill=SoftLavender,right=7mm of c] (d) {\textbf{8--10 minutes}\\Choose one action\\Explain why it supports belonging};
\draw[arrow] (a)--(b); \draw[arrow] (b)--(c); \draw[arrow] (c)--(d);
\end{tikzpicture}}
\caption{The 10-minute lesson sequence. The visual timeline reduces uncertainty and connects each teaching move to its purpose.}
\label{fig:lesson-timeline}
\end{figure}

I explained that autism is one way a person's brain can work and that autistic people are not all the same. Communication may include speech, gestures, writing, pictures, AAC, movement, or silence. Some students may need quiet, space, clear directions, processing time, predictable routines, or a break from noise. I did not ask students to identify anyone's diagnosis or name an autistic person they knew.

Using the OAR friendship guidance, I taught five actions:
\begin{enumerate}
\item \textbf{Invite} a classmate without forcing participation.
\item \textbf{Wait} so the person has time to process, communicate, or try.
\item \textbf{Ask} respectfully instead of guessing what someone needs.
\item \textbf{Respect} sensory needs, personal space, communication tools, breaks, and privacy.
\item \textbf{Include} students who are new, quiet, tired, or participating differently as full group members.
\end{enumerate}

The main practice scenario came from placement patterns: during a coding activity, one student already knew how to enter the program, several classmates were still stuck, and another learner was new to the group. Students suggested showing one step without grabbing the Chromebook, using a calm voice, waiting for the learner to try, and treating the new learner as part of the class. The key distinction was between helping and taking over. The device and decision-making remained with the learner.

I closed with: \textit{One way I can be a supportive classmate is... because...} The word ``because'' required students to connect an action to reducing stress, increasing access, protecting dignity, or helping someone belong.
